\chapter{Discussion}

\section{Synthesis: three separable
problems}

The clearest lesson of Chapter 5 is that the end-to-end failures of a
video RAG system decompose into three problems that are \emph{separable}
--- they have different causes, respond to different interventions, and
the interventions compose.

The first is a \textbf{recall problem}: the correct moment never enters
the candidate pool because the query is phrased in causal or temporal
language that a contrastive vision-language encoder cannot match against
pixels. Query decomposition addresses exactly this --- rewriting the
question into scene captions lifted deep-rank recall (corpus R@10 +10\%
relative) while leaving the top rank untouched (§5.3.2). The second is a
\textbf{precision problem}: the correct moment is in the pool but
near-tied cosine scores cannot order it first. Second-stage re-scoring
with a stronger visual model addresses this, and only this (§5.3.3); so
does exchanging the index backbone outright (§5.3.4). The third is an
\textbf{evidence problem}: the correct moment is retrieved, correctly
time-stamped, and placed in front of the answerer --- which then cannot
answer because the decisive evidence is visual and its input is text. No
amount of retrieval improvement fixes this; changing the evidence
channel does (+87\% on temporal-next questions, §5.4.2).

Reading the results this way explains why the interventions compose
additively rather than cannibalising each other (Table 5.2): each
removes a different bottleneck. It also offers a diagnostic recipe for
practitioners: compare corpus against video scope to separate
cross-video confusion from localisation error, compare shallow against
deep recall to separate recall from precision problems, and compare
text-only against multimodal answering to expose evidence starvation.

A second thread worth drawing out is that \textbf{the value of a
modality is role- and question-dependent, not absolute}. The transcript
channel failed twice as a \emph{retrieval key} --- as embeddings (§5.2)
and as a cross-encoder signal (§5.3.3) --- and the prior-only baselines
of §5.4.3 discipline any remaining optimism about its third role as
\emph{evidence}: the best text-only stack (.547) merely matches the bare
model's answer prior (.560), and the apparent causal gains largely
dissolve against that bar (CW .590 against a .574 prior). What survives
is sharply type-dependent. Transcript evidence genuinely beats the prior
where speech describes the concurrent scene (TC .773 vs .545) and falls
far below it where the answer is a silent visual action (TN .341 vs
.568) --- an instructed-grounding effect in which the model abandons a
useful prior in favour of evidence that contains nothing. Speech in
everyday video is a poor pointer to \emph{where} something happens, and
describes \emph{what} happens only when someone happens to narrate it.
Modality ablations that report a single ``contribution'' number conflate
roles and question types; our results suggest both splits are necessary
--- and that any evidence-grounded system should be benchmarked against
the bare model's prior, not against random.

Finally, the two-stage equivalence result (§5.3.4) has an architectural
implication beyond this project. At 5,725 segments, indexing with the
expensive model is trivially affordable, and yet the
cheap-index/expensive-re-rank design already matched it to within
measurement noise. Since re-ranking cost is fixed (top-30 per query)
while indexing cost grows linearly with the corpus, the two-stage design
is the only one of the two that survives scaling --- and this work
provides evidence that the quality sacrifice can be nil.

\section{Limitations}

\textbf{Domain coverage.} The primary conclusions are drawn from
NExT-QA/GQA: short, everyday, speech-sparse home videos. The negative
results about the text channel --- fusion never helping (§5.3.1), the
text cross-encoder hurting (§5.3.3) --- are \emph{conditional on this
domain}, and that conditionality is now a measured result rather than a
caveat: the speech-dense replication of §5.5 (YouCook2, 40 videos)
reversed both signs exactly as predicted, with transcript retrieval
doubling visual retrieval and late fusion becoming the best
configuration. What remains limited is coverage: two domains, a small
second corpus, and a retrieval-only replication --- the positive results
(decomposition, visual re-ranking, backbone ordering, evidence-channel
effect) have not yet been re-tested outside NExT-QA, and the QA-stage
experiments were not repeated on the second domain. Their transfer
argument still rests on properties of contrastive encoders and of the QA
task rather than on measurement.

\textbf{A confound inside the text-channel result.} The weak text
channel has two entangled causes: the transcripts are often
\emph{uninformative} (chatter unrelated to the questions), and they are
often \emph{unreadable} by the predominantly-English CLIP text encoder
(multilingual speech). The present experiments cannot apportion blame
between the two. A translate-then-embed variant --- machine-translating
all transcripts to English before indexing --- would separate them
cheaply and was left undone only for time.

\textbf{Segmentation granularity bounds temporal accuracy.} Segments are
fixed 8-second windows on a 4-second stride, while ground-truth
intervals vary freely; even a perfect ranker cannot exceed the tIoU that
the best-aligned window achieves. Part of the headline gap between τ =
0.5 and τ = 0.3 results (§5.3.5) is this quantisation, not ranking
error. tIoU@1 therefore carries a structural ceiling that is a property
of the segmenter; adaptive or shot-aligned segmentation would raise the
ceiling itself.

\textbf{Multiple-choice accuracy is a proxy.} NExT-QA's five-way format
enables cheap, objective scoring, but it measures answer
\emph{selection}, not answer \emph{quality}: the faithfulness of the
generated prose and the correctness of its timestamp citations are
audited only by the agent's own reflection pass, not by human or
LLM-as-judge evaluation. The white-dog case study (§5.4.2) illustrates
that the generated answers can be precise and well-grounded, but no
systematic open-ended-answer evaluation was performed.

\textbf{Sample sizes and single runs.} Retrieval results cover all 3,358
questions and are deterministic given the frozen decomposition cache.
The QA experiments, however, are budget-limited: 150 questions for the
agent comparison, 44 for the controlled vision experiment, and the text
cross-encoder was rejected on a 100-question development subset. All
agent numbers are single runs of non-deterministic API models. The two
central paired comparisons are, however, statistically supported: exact
McNemar tests on the shared question sets give p = .031 for the agent
effect and p = .004 for the evidence-channel effect, with
paired-bootstrap 95\% CIs of {[}+.013, +.187{]} and {[}+.114, +.455{]}
respectively (§5.4.1, §5.4.2) --- so neither headline effect is
plausibly run-to-run variance. The caution that remains is for the
\emph{small} differences in Table 5.5's by-type columns (n ≤ 22), which
carry no such support and should not be over-read.

\textbf{Asymmetries in the agent comparison.} The controlled vision
experiment (§5.4.2) ran under the \emph{simple} agent, because at that
point the LangGraph harness targeted OpenAI-compatible providers only;
the graph-plus-vision combination was implemented subsequently and
evaluated in §5.4.4, confirming the floor-not-ceiling reading on the
overall figure (.727 under the combined configuration). Its TN figure,
however, is .568 --- three questions \emph{below} the simple agent's
.636 on the same items. That gap is within noise at n = 44, but it
points at the asymmetry that remains: whether the reflection pass
benefits from \emph{seeing} the keyframes it audits, rather than the
text rendering of the evidence it currently audits, has not been
isolated, and TN is precisely the type where a text-mediated audit has
least to offer.

\textbf{No user study.} The project plan scheduled a small user study
(W10); it was descoped in favour of the controlled vision experiment and
the ablation matrix, which we judged to produce more defensible evidence
per unit time. The Streamlit demo (§3.6) stands in for qualitative
validation, but claims about end-user utility remain unsupported and are
deliberately absent from this dissertation.

\section{Threats to validity}

\textbf{Construct validity.} The hit criterion (correct video ∧ tIoU ≥
τ) is stricter than video-level retrieval metrics common in prior work;
this deflates absolute numbers relative to that literature but is
faithful to the moment-retrieval construct. For QA, the multiple-choice
format admits answer-prior effects: a language model can reject
implausible distractors without evidence. The controlled experiment
bounds this --- priors were identical across both evidence conditions,
so the +29.5-point effect cannot be prior-driven --- and the prior-only
baselines (§5.4.3) measure the prior directly for both models (.560 each
with no evidence at all), completing the decomposition: the .341
text-only TN figure turned out to sit \emph{below} the model's own
prior, so what looked like weak evidence was evidence-induced
suppression. This is why the dissertation rests its QA claims on paired
within-model comparisons rather than absolute accuracies.

\textbf{Internal validity.} The decomposition cache is shared verbatim
across all configurations that use it, eliminating decomposition
variance from those comparisons. The re-ranker's fusion weight (0.5) and
candidate pool (30) were fixed a priori rather than tuned, and the
development subset used to reject the text cross-encoder (n = 100) is
disjoint in purpose but not in items from the full evaluation set; since
the decision it informed was to \emph{exclude} a component, any bias is
conservative. Agent experiments called the same LLM family for
decomposition and answering (DeepSeek), which could in principle
correlate their errors; retrieval experiments are unaffected, as
decomposition quality is identical across all retrieval configurations.

\textbf{External validity.} Beyond the domain limitation above, two
scale caveats apply. The corpus (567 videos, 5,725 segments) is small
enough that exhaustive HNSW search is effectively exact; ANN recall
degradation at millions of segments is not measured. And the two-stage
equivalence claim is demonstrated at a scale where \emph{both} designs
are feasible --- its extrapolation to large corpora rests on the
argument of §6.1, not on direct measurement. Finally, all API-based
results are tied to unversioned commercial model endpoints
(deepseek-chat, claude-opus-4-8, mid-2026); the committed prompts,
caches, and result files make the analysis auditable, but bit-exact
replication of agent runs is not guaranteed by the providers.

\section{Summary}

The findings hold together as a single narrative: in everyday video,
visual signal dominates at every stage where a modality choice exists
--- and the speech-dense replication shows this dominance is a
measurable property of the domain, reversing cleanly where narration is
dense and task-aligned. The remaining errors are separable into recall,
precision, and evidence problems with independent, composable remedies,
benchmarked against the bare model's prior rather than against random.
The principal limitations --- partial domain coverage, proxy QA metrics,
budget-limited agent samples, and the untested multimodal-reflection
interaction --- are all addressable with more compute rather than new
machinery, and none threatens the direction of the reported effects,
whose magnitudes (+69\% relative R@1 from the backbone, +87\% from the
evidence channel, 2× modality reversal on YouCook2) far exceed the
plausible noise floor. Chapter 7 concludes and sets out the
corresponding future work.
